\chapter{Building Block View (Static View)}

The SUM_TCS module is a SWC within the context of AUTOSAR.
It consists of static files and generated files.

\section{Configuration}

SUM_TCS configuration is done the same way as for standard BSW modules.
A parameter definition file is included with the module which allows the users to configure the available parameters.
The generated files change based on the configuration.
The configuration parameters for TCS are provided in the paramdef file \file{SUM_TCS_EcuParamDef.arxml}.
The user shall configure values to the various parameters provided in the paramdef: based on the inputs configured, TCS shall use OAW scripts to generate the source and header files along with the SWCD.

\subsection{Included Containers:}

\begin{longtable}{|l|l|l|}
\hline
\begin{minipage}{0.18\linewidth}\raggedright\vspace{.75ex}\textbf{Name}\vspace{.75ex}\end{minipage} & \begin{minipage}{0.18\linewidth}\raggedright\vspace{.75ex}\textbf{Multiplicity}\vspace{.75ex}\end{minipage} & \begin{minipage}{0.55\linewidth}\raggedright\vspace{.75ex}\textbf{Scope / Dependency}\vspace{.75ex}\end{minipage}\\
\hline\hline
\endfirsthead
\hline
\begin{minipage}{0.18\linewidth}\raggedright\vspace{.75ex}\textbf{Name}\vspace{.75ex}\end{minipage} & \begin{minipage}{0.18\linewidth}\raggedright\vspace{.75ex}\textbf{Multiplicity}\vspace{.75ex}\end{minipage} & \begin{minipage}{0.55\linewidth}\raggedright\vspace{.75ex}\textbf{Scope / Dependency}\vspace{.75ex}\end{minipage}\\
\hline\hline
\endhead
\hline
\endfoot
\hline
\endlastfoot
\begin{minipage}{0.18\linewidth}\raggedright\vspace{.75ex}TcsGeneral\vspace{.75ex}\end{minipage} & \begin{minipage}{0.18\linewidth}\raggedright\vspace{.75ex}1..1\vspace{.75ex}\end{minipage} & \begin{minipage}{0.55\linewidth}\raggedright\vspace{.75ex}General configuration paramers for TCS\vspace{.75ex}\end{minipage}\\\hline
\end{longtable}

\subsection{TcsGeneral}
\label{sec:TcsGeneral}


\Requirement{031}


{
\ConfigurationParameter{MasterClockPollingPeriodicity}
{Time in msec for the periodicity of MasterClockPolling}
{1..1}
{ECUC-INTEGER-PARAM-DEF}
{1..1000}
{100}

\ConfigurationParameter{PermitExecutionMask}
{Bitmask used to identify the Wakeup bit in the \code{SUSD_TCS_PermitExecution} interface.}
{1..1}
{ECUC-INTEGER-PARAM-DEF}
{0..255}
{254}

\ConfigurationParameter{SumTCSGlobalTimePrecision}
{Local time source accuracy in msec. Must be a positive multiple of 100 ms and be greater than or equal to TimeManagementPeriodicity.}
{1..1}
{ECUC-INTEGER-PARAM-DEF}
{1..1000}
{100}

\ConfigurationParameter{TimeManagementPeriodicity}
{Time in msec for the periodicity of TimeManagement}
{1..1}
{ECUC-INTEGER-PARAM-DEF}
{1..1000}
{100}

\ConfigurationParameter{TcsGenerateInitRunnable}
{Enable or Disable the Runnable Init feature.}
{1..1}
{ECUC-BOOLEAN-PARAM-DEF}
{NA}
{false}

\ConfigurationParameter{TcsInitFuncPeriodicity}
{Specifies time in ms for the periodicity of the runnable Runnable_TCS_Init.}
{0..1}
{ECUC-FLOAT-PARAM-DEF}
{1.0..5000.0}
{1000.0}

\ConfigurationParameter{TCS\_TimeSyncMessageSignalRef}
{Reference to the ComSignal or ComGroupSignal used for the TimeSyncMessage reception. Determines the signal type used for master clock polling.}
{1..1}
{ECUC-CHOICE-REFERENCE-DEF}
{-..-}
{-}

\ConfigurationParameter{TCS\_ServiceBased\_TimeSyncMessage\_Ref}
{Reference to the implementation data type element for service-based TimeSyncMessage communication. When configured, TCS uses SOME/IP service-based reception instead of signal-based.}
{0..1}
{ECUC-FOREIGN-REFERENCE-DEF}
{-..-}
{-}

\ConfigurationParameter{TCS\_ServiceBased\_TimeSyncMessage\_DataTypeMappingSet}
{Reference to the data type mapping set required for service-based TimeSyncMessage communication. Must be configured when \code{TCS\_ServiceBased\_TimeSyncMessage\_Ref} is set.}
{0..1}
{ECUC-FOREIGN-REFERENCE-DEF}
{-..-}
{-}

\ConfigurationParameter{DID_F129_TotalVehicleTimeValue_Reference}
{Reference to DCM DID 0xF129 Total Vehicle Time Value}
{1..1}
{ECUC-REFERENCE-DEF}
{-..-}
{-}

\ConfigurationParameter{TcsDemTimestampRef}
{Reference to Dem Timestamp}
{1..1}
{ECUC-REFERENCE-DEF}
{-..-}
{-}
}
{\Satisfies{SRS_TCS 019}}
\section{Calibration Parameters}
N/A
